\section{Experienced pilots and beginners}\label{sec:levels}

Does the route explain the same share for both wing levels?

It does (Table~\ref{tab:isolated}, Figure~\ref{fig:isolated}). For the \IsoPiecesExp{} pieces of
experienced pilots $\Delta\Hloc(1000~\mathrm{s})$ falls from \IsoBaseDHlocExp{} to
\IsoDHlocExp{} \IsoDHlocExpBand, \IsoExplainedExp\% explained. For the \IsoPiecesBeg{} pieces of
beginners it falls from \IsoBaseDHlocBeg{} to \IsoDHlocBeg{} \IsoDHlocBegBand, \IsoExplainedBeg\%
explained. The beginners' baseline gap is larger and their band wider, since they have fewer pieces;
the two levels differ by one point in the share explained, well within their bands.

\begin{figure}
\centering
\includegraphics[width=\textwidth]{isolated.pdf}
\caption{Left: $\Hloc(\tau)$ of group time (orange) and of the same route flown alone (blue), by wing
level; dashed, the long series. Right: the difference group minus control of two $\Hloc$ and of the
statistics of climbs and glides, with 5--95\% bands, by wing level.}
\label{fig:isolated}
\end{figure}

\paragraph{What we keep.} Experienced pilots and beginners give the same share explained, so that the
route accounts for the same part of the gap whatever the wing level.
